\documentclass[a4paper]{article}

\newcommand{\docTitle}{Robotik}

\usepackage{datetime2}  % for dates
\usepackage[utf8]{inputenc} % UTF-8 input encoding
\usepackage[T1]{fontenc} % fixes \textquotedbl with pdflatex/OT1
\usepackage{textcomp}   % additional text symbols for typewriter/upquote
\usepackage{xcolor}     % for gray color
\usepackage{listings}
\usepackage{marginnote} % for horizontal margin notes
\usepackage{parskip}
\usepackage{amsmath,amssymb,mathtools}
\usepackage{everypage}
\usepackage{fancyhdr}
\usepackage{amsfonts}   % Mengen and natural numbers
\usepackage{enumerate}  % Custom Enumerations
\usepackage[inline]{enumitem}   % List spacing and labels
\usepackage{multicol}   % Multiple Columns
\usepackage{tabularx}

\usepackage{hyperref}   % Links
\hypersetup{
    pdftitle={\docTitle},
    pdfauthor={Moritz Köhler},
    pdfkeywords={DHBW, Hochschule},
    colorlinks=true,
    linkcolor=black,
    filecolor=magenta,
    urlcolor=cyan
}

% -------
% Images
% -------

\usepackage{import}
\usepackage{xifthen}
\usepackage{pdfpages}
\usepackage{tikz}
\usetikzlibrary{arrows.meta, positioning}

\newcommand{\incfig}[1]{%
    \def\svgwidth{\columnwidth}
    \import{./fig/}{#1.pdf_tex}
}

% ---------------------------------
% Vertical Side note (Bottom Left)
% ---------------------------------

\newcommand{\verticaldocinfo}{%
  \rotatebox{90}{\tiny\color{gray}\texttt{\DTMtoday\_\docTitle\_\jobname}}%
}

\AddEverypageHook{
  \put(-40pt,-650pt){\verticaldocinfo}%
}

% ---------------
% Math Shortcuts
% ---------------

% Number sets
\newcommand{\R}{\mathbb{R}}
\newcommand{\N}{\mathbb{N}}
\newcommand{\Z}{\mathbb{Z}}
\newcommand{\Q}{\mathbb{Q}}
\newcommand{\C}{\mathbb{C}}

% Vectors and matrices
\newcommand{\vect}[1]{\mathbf{#1}}
\newcommand{\mat}[1]{\mathbf{#1}}

% Derivatives
\newcommand{\dd}{\mathrm{d}}
\newcommand{\dpartial}[2]{\frac{\partial #1}{\partial #2}}
\newcommand{\ddt}{\frac{\dd}{\dd t}}

% Norms and absolute values
\newcommand{\norm}[1]{\left\lVert #1 \right\rVert}
\newcommand{\abs}[1]{\left\lvert #1 \right\rvert}

% Probability & Statistics
\newcommand{\E}{\mathbb{E}}
\newcommand{\Var}{\mathrm{Var}}
\newcommand{\Prob}{\mathbb{P}}

% Fields and Algebras
\newcommand{\K}{\mathbb{K}}
\newcommand{\F}{\mathbb{F}}

% Set Operations
\newcommand{\compl}[1]{#1^{\mathrm{c}}}

% Matrix and Vector Operations
\newcommand{\T}{^{\top}}
\newcommand{\inv}{^{-1}}
\newcommand{\conj}[1]{\overline{#1}}

% -------------------
% Main lecture macro 
% -------------------

\newcommand{\lecture}[1]{%
  \protect\marginnote{\protect\tiny\protect\color{gray}#1}
}

% -----------------------------
% Command for box with heading
% -----------------------------

\fboxsep2mm
\newcommand{\know}[2]{
    \noindent\fbox{
        \parbox{\textwidth-21pt}{

            \textbf{#1:}
            \vspace{0.125cm}

            #2
        }
    }
}

% -------------------
% Listings Languages
% -------------------

% Shared defaults for code listings
\lstset{
    basicstyle=\ttfamily\small,
    keywordstyle=\color{blue}\bfseries,
    commentstyle=\color{gray},
    stringstyle=\color{teal},
    numberstyle=\tiny\color{gray},
    numbers=left,
    stepnumber=1,
    numbersep=8pt,
    showstringspaces=false,
    frame=single,
    breaklines=true,
    tabsize=4,
    keepspaces=false,
    columns=flexible,
    upquote=true,
    extendedchars=false % disable extended char handling to avoid UTF-8 conflicts
}

% SQL syntax highlighting
\lstdefinelanguage{SQL}{
    morekeywords={SELECT,FROM,WHERE,INSERT,INTO,VALUES,UPDATE,SET,DELETE,CREATE,TABLE,PRIMARY,KEY,NOT,NULL,AND,OR,JOIN,LEFT,RIGHT,INNER,ON,AS,ORDER,BY,GROUP,HAVING,LIMIT},
    sensitive=false,
    morecomment=[l]{--},
    morestring=[b]'
}

% JSON syntax highlighting
\lstdefinelanguage{json}{
    morestring=[b]",
    stringstyle=\color{teal},
    comment=[l]{//},
    morecomment=[s]{/*}{*/},
        morekeywords={true,false,null},
    literate=
     *{0}{{{\color{blue}0}}}{1}
      {1}{{{\color{blue}1}}}{1}
      {2}{{{\color{blue}2}}}{1}
      {3}{{{\color{blue}3}}}{1}
      {4}{{{\color{blue}4}}}{1}
      {5}{{{\color{blue}5}}}{1}
      {6}{{{\color{blue}6}}}{1}
      {7}{{{\color{blue}7}}}{1}
      {8}{{{\color{blue}8}}}{1}
      {9}{{{\color{blue}9}}}{1}
            {:}{{{\color{black}:}}}{1}
            {,}{{{\color{black},}}}{1},
    showstringspaces=false
}

% Language specific styles
\lstdefinestyle{javastyle}{
    language=Java,
    basicstyle=\ttfamily\small,
    keywordstyle=\color{blue}\bfseries,
    commentstyle=\color{gray},
    stringstyle=\color{teal},
    numberstyle=\tiny\color{gray},
    numbers=left,
    stepnumber=1,
    numbersep=8pt,
    showstringspaces=false,
    frame=single,
    breaklines=true,
    tabsize=4,
    keepspaces=true,
    columns=fullflexible,
    upquote=true
}

\lstdefinestyle{sqlstyle}{
    language=SQL,
    basicstyle=\ttfamily\small,
    keywordstyle=\color{blue}\bfseries,
    commentstyle=\color{gray},
    stringstyle=\color{teal},
    numberstyle=\tiny\color{gray},
    numbers=left,
    stepnumber=1,
    numbersep=8pt,
    showstringspaces=false,
    frame=single,
    breaklines=true,
    tabsize=4,
    keepspaces=true,
    columns=fullflexible,
    upquote=true
}

\lstdefinestyle{pythonstyle}{
    language=Python,
    basicstyle=\ttfamily\small,
    keywordstyle=\color{blue}\bfseries,
    commentstyle=\color{gray},
    stringstyle=\color{teal},
    numberstyle=\tiny\color{gray},
    numbers=left,
    stepnumber=1,
    numbersep=8pt,
    showstringspaces=false,
    frame=single,
    breaklines=true,
    tabsize=4,
    keepspaces=true,
    columns=fullflexible,
    upquote=true
}

\lstdefinestyle{cstyle}{
    language=C,
    basicstyle=\ttfamily\small,
    keywordstyle=\color{blue}\bfseries,
    commentstyle=\color{gray},
    stringstyle=\color{teal},
    numberstyle=\tiny\color{gray},
    numbers=left,
    stepnumber=1,
    numbersep=8pt,
    showstringspaces=false,
    frame=single,
    breaklines=true,
    tabsize=4,
    keepspaces=true,
    columns=fullflexible,
    upquote=true
}

\lstdefinestyle{bashstyle}{
    language=bash,
    basicstyle=\ttfamily\small,
    keywordstyle=\color{blue}\bfseries,
    commentstyle=\color{gray},
    stringstyle=\color{teal},
    numberstyle=\tiny\color{gray},
    numbers=left,
    stepnumber=1,
    numbersep=8pt,
    showstringspaces=false,
    frame=single,
    breaklines=true,
    tabsize=4,
    keepspaces=true,
    columns=fullflexible,
    upquote=true
}

\lstdefinestyle{jsonstyle}{
    language=json,
    basicstyle=\ttfamily\small,
    keywordstyle=\color{blue}\bfseries,
    commentstyle=\color{gray},
    stringstyle=\color{teal},
    numberstyle=\tiny\color{gray},
    numbers=left,
    stepnumber=1,
    numbersep=8pt,
    showstringspaces=false,
    frame=single,
    breaklines=true,
    tabsize=4,
    keepspaces=true,
    columns=fullflexible,
    upquote=true
}

% Default listing style
\lstset{language={}}

\title{\docTitle}
\author{Moritz}
\date{\today}

\begin{document}
\maketitle
\tableofcontents

\lecture{2026-12-12}

\section{Modellierung und Simulation}

\subsection{Kinematik}

Wir brauchen Position und Drehung eines Punktes. Also 6 Freiheitsgrade.

Homogene Transformationsmatrizen

Wir können mit der Direkte (Vorwärts) Kinematik aus den Winkeln und den Längen der Armen den Endpunkt berechnen.

Die Invers Kinematik ist dann die Mathematische Abbildung von einem Ziel (Punkt) die Gelenkwinkel zu bestimmen.

\know{Wichtig}{Folie 83 ist besonders Prüfungsrelevant. Zudem sollten die Formel von S. 84}

\begin{align*}
    x^2_{TCP} + y^2_{TCP} &= l_1^2+l_2^2 - 2 l_1l_2\cos(\pi - q_2)\\
    q_2 &= \pi - \arccos\left(\frac{l_1^2+l_2^2-x_{TCP}^2-y_{TCP}^2}{2l_1l_2}\right)
\end{align*}

\begin{align*}
    \beta &= \arctan\left(\frac{y_{TCP}}{x_{TCP}}\right)\\
    q_1 &= \beta + \gamma
\end{align*}

$$\gamma = \arccos\left(\frac{l_1^2+x_{TCP}^2+y_{TCP}^2-l_2^2}{2l_1\sqrt{x_{TCP}^2+y_{TCP}^2}}\right)$$

Es gibt also zwei Räume:

\begin{enumerate}
    \item Gelenkwinkelraum
    \item Kartesischer Raum
\end{enumerate}

\subsection{Optimierung}

Jaccobi-Matrix und Gradientenverfahren

\subsection{Differentielle Kinematik}

Gibt es auch noch in invers

DH Transformation:

\begin{enumerate}
    \item Rotation $\theta_n$
    \item Translation $d_n$
    \item Translation $a_n$
    \item Rotation $\alpha_n$
\end{enumerate}

\subsection{Dynamik}

\begin{enumerate}
    \item Lagrange-Formalismus (Einfach, Systematisch, eignet sich für händisch)
    \item Newton-Euler-Formalismus (Rekursiv)
\end{enumerate}

\subsection{Lagrange-Formalismus}

Kochrezept dafür (S. 99):

\begin{enumerate}
    \item Festelgen N verallgemeinerten koordinaten $q_j$
    \item Aufstellen der Ortsvektoren zu den Massen: $r_j = f(q, t), v_j = \frac{dr_j}{dt}$
    \item Aufstellen der kinetischen Energie $E_{kin}$ un der potentiellen Energie $E_{pot}$
    \item Aufstellen der Lagrange-Funktion: $L=E_{kin} - E_{pot}$
    \item Ermittlung der Euler-Lagrange-Gleichung für alle $j=1, ..., N$: $$\frac{d}{dt}\left(\frac{\delta L}{\delta \dot q_j}\right)-\frac{\delta L}{\delta q_j} = Q_j$$
\end{enumerate}

\subsection{Simulation}

Es gibt mehrere Entwicklungswerkzeuge für die Robotik:

\begin{itemize}
    \item Matlab/ Simulink
    \item ROS/ ROS2
    \item Nvidia Isaac Sim
\end{itemize}

\subsubsection{Rapid Control Prototyping}

Es geht um die Methodik:

\begin{enumerate}
    \item Beschreibung und Modellierung des zu automatisierenden (dynamischen) Systems
    \item Regelungs- und Steuerungsentwurf anhand des abgeleiteten Modells
    \item Erprobung der entwickelten Algorithmen in einer Simulationsumgebung
    \item Umsetzung der entwickelten Algorithmen auf dem Steuergerät (Target)
    \item Erprobung der entwickelten (in Simulation validierten) Algorithmen am realen System
\end{enumerate}

Hil (Hardware in the Loop), Sil (System in the Loop, Code Generierung), MiL (Model in the Loop)

\subsection{Matlab/ Simulink}

\href{https://de.mathworks.com/campaigns/offers/matlab-basic-functions-reference-cheat-sheet.html}{Cheat Sheet}

\end{document}